\chapter{1977年西藏卷}

\section{解答题}

\begin{problem}
计算：\(\left(4.8-3.7\right)^2-\sqrt[3]{-27}\times\left(5-1\div5\right)\times\left(-\frac{5}{8}\right)\).
\end{problem}

\begin{solution}
先分别计算各部分：\(4.8-3.7=1.1\)，\(\sqrt[3]{-27}=-3\)，且 \(5-1\div5=5-\frac{1}{5}=\frac{24}{5}\). 因此原式为
\[
1.1^2-\left(-3\right)\times\frac{24}{5}\times\left(-\frac{5}{8}\right)=1.21-9=-7.79
\]
故原式的值为 \(-7.79\).
\end{solution}

\begin{problem}
求分式方程 \(\frac{1-x}{x}=x\) 在 \(0\) 到 \(1\) 范围内精确到小数三位的根.
\end{problem}

\begin{solution}
原方程的分母为 \(x\)，所以定义域要求 \(x\ne0\). 在此条件下，两边同乘 \(x\)，得 \(1-x=x^2\)，即 \(x^2+x-1=0\). 由求根公式，
\[
x=\frac{-1\pm\sqrt{1+4}}{2}=\frac{-1\pm\sqrt{5}}{2}
\]
因为这两个根的乘积为 \(-1\)，所以它们均不为零；又因为去分母是在 \(x\ne0\) 的定义域内进行的，所以它们都是原分式方程的根. 由 \(2<\sqrt{5}<3\) 可知 \(\frac{-1-\sqrt{5}}{2}<0\)，且 \(0<\frac{-1+\sqrt{5}}{2}<1\).
故 \(0\) 到 \(1\) 范围内的根为
\[
x=\frac{-1+\sqrt{5}}{2}\approx\frac{-1+2.236}{2}=0.618
\]
即精确到小数三位为 \(0.618\).
\end{solution}

\begin{problem}
已知 \(\log_{10}2=0.3010\)，\(\log_{10}5=0.6990\)，试计算 \(\left(2\log_{10}4\right)\left(\log_2 5\right)+\log_{10}\frac{5}{2}\) 的值（精确到小数三位）.
\end{problem}

\begin{solution}
由对数的幂运算法则、换底公式和商的对数公式，分别有 \(\log_{10}4=2\log_{10}2\)，\(\log_2 5=\frac{\log_{10}5}{\log_{10}2}\)，\(\log_{10}\frac{5}{2}=\log_{10}5-\log_{10}2\). 所以
\[
\begin{gathered}
\left(2\log_{10}4\right)\left(\log_2 5\right)+\log_{10}\frac{5}{2}\\
=\left(4\log_{10}2\right)\frac{\log_{10}5}{\log_{10}2}+\log_{10}5-\log_{10}2\\
=4\log_{10}5+\log_{10}5-\log_{10}2=5\log_{10}5-\log_{10}2\\
=5\times0.6990-0.3010=3.1940
\end{gathered}
\]
因此精确到小数三位为 \(3.194\).
\end{solution}

\begin{problem}
利用特殊角三角比，求 \(\sin75^\circ\) 的值.
\end{problem}

\begin{solution}
将 \(75^\circ\) 写成 \(45^\circ+30^\circ\)，利用正弦的和角公式，得
\[
\sin75^\circ=\sin\left(45^\circ+30^\circ\right)=\sin45^\circ\cos30^\circ+\cos45^\circ\sin30^\circ
\]
代入特殊角的三角比，得
\[
\sin75^\circ=\frac{\sqrt{2}}{2}\cdot\frac{\sqrt{3}}{2}+\frac{\sqrt{2}}{2}\cdot\frac{1}{2}=\frac{\sqrt{6}+\sqrt{2}}{4}
\]
故 \(\sin75^\circ=\frac{\sqrt{6}+\sqrt{2}}{4}\).
\end{solution}

\begin{problem}
证明恒等式：\(\left(\cos\alpha-\cos\beta\right)^2+\left(\sin\alpha-\sin\beta\right)^2=2-2\cos\left(\alpha-\beta\right)\).
\end{problem}

\begin{solution}
展开左端，并利用 \(\sin^2\theta+\cos^2\theta=1\)，得
\[
\begin{gathered}
\left(\cos\alpha-\cos\beta\right)^2+\left(\sin\alpha-\sin\beta\right)^2\\
=\cos^2\alpha-2\cos\alpha\cos\beta+\cos^2\beta+\sin^2\alpha-2\sin\alpha\sin\beta+\sin^2\beta\\
=2-2\left(\cos\alpha\cos\beta+\sin\alpha\sin\beta\right)
\end{gathered}
\]
再由余弦的差角公式 \(\cos\left(\alpha-\beta\right)=\cos\alpha\cos\beta+\sin\alpha\sin\beta\)，可得
\[
\left(\cos\alpha-\cos\beta\right)^2+\left(\sin\alpha-\sin\beta\right)^2=2-2\cos\left(\alpha-\beta\right)
\]
因此恒等式成立.
\end{solution}

\section{选做题}

\begin{problem}
红旗工厂要围成一个周长为 \(120\text{ 米}\) 的矩形场地，可以存放原料. 问矩形的长和宽应各取多少米，才能使存放场地的面积最大？并说明这最大面积的场地是什么形状？
\end{problem}

\begin{solution}
设矩形的长为 \(x\text{ 米}\)，宽为 \(y\text{ 米}\). 由周长条件得 \(2x+2y=120\)，所以 \(y=60-x\). 由于长、宽均为正数，必须有 \(0<x<60\). 面积为
\[
S=xy=x\left(60-x\right)=-x^2+60x=-\left(x-30\right)^2+900
\]
因为 \(\left(x-30\right)^2\geq0\)，所以 \(S\leq900\)，且当且仅当 \(x=30\) 时取等号. 此时 \(y=60-30=30\). 故矩形的长、宽均取 \(30\text{ 米}\) 时面积最大，最大面积为 \(900\text{ 平方米}\)；此时长和宽相等，所以最大面积的场地是正方形.
\end{solution}

\begin{problem}
有一正四棱台，上底边长是 \(40\text{ 厘米}\)，下底边长是 \(52\text{ 厘米}\)，高是 \(8\text{ 厘米}\). 求：

\begin{enumerate}
\item 它的下底面的面积是多少？
\item 它的侧面积是多少？
\item 它的容积是多少？
\end{enumerate}
\end{problem}

\begin{solution}
\begin{enumerate}
\item 正四棱台的下底面是边长为 \(52\text{ 厘米}\) 的正方形，因此下底面面积为 \(S_{\text{下底}}=52^2=2704\text{ 平方厘米}\).

\item 正四棱台的每个侧面都是全等的等腰梯形，且梯形的上、下底分别为 \(40\text{ 厘米}\) 和 \(52\text{ 厘米}\). 过梯形上、下底的中点作垂线，底边差的一半为 \(\frac{52-40}{2}=6\text{ 厘米}\). 因此梯形的高，即正四棱台的斜高为
\[
l=\sqrt{8^2+6^2}=10\text{ 厘米}
\]
四个侧面均为等腰梯形，所以侧面积为
\[
S_{\text{侧}}=4\times\frac{40+52}{2}\times10=1840\text{ 平方厘米}
\]

\item 上、下底面的面积分别为 \(S_1=40^2=1600\text{ 平方厘米}\)，\(S_2=52^2=2704\text{ 平方厘米}\). 棱台的体积公式为 \(V=\frac{h}{3}\left(S_1+S_2+\sqrt{S_1S_2}\right)\). 由 \(\sqrt{S_1S_2}=\sqrt{40^2\times52^2}=40\times52=2080\)，且 \(h=8\text{ 厘米}\)，得
\[
V=\frac{8}{3}\left(1600+2704+2080\right)=\frac{8}{3}\times6384=17024\text{ 立方厘米}
\]
\end{enumerate}
\end{solution}

\begin{problem}
已知直线 \(L_2\) 经过点 \(\left(7,1\right)\)，并且与直线 \(L_1:y=\frac{1}{3}x+2\) 垂直，求直线 \(L_2\) 的方程.
\end{problem}

\begin{solution}
直线 \(L_1\) 的斜率为 \(\frac{1}{3}\). 设直线 \(L_2\) 的斜率为 \(m\). 两条非竖直直线垂直时斜率之积为 \(-1\)，所以 \(m\cdot\frac{1}{3}=-1\)，即 \(m=-3\). 利用点斜式并代入点 \(\left(7,1\right)\)，得 \(y-1=-3\left(x-7\right)\)，整理得 \(y=-3x+22\). 因此直线 \(L_2\) 的方程为 \(y=-3x+22\).
\end{solution}

\section{参考题}

\begin{problem}
求证：\(1^3+2^3+3^3+\cdots+n^3=\left[\frac{n\left(n+1\right)}{2}\right]^2\).
\end{problem}

\begin{solution}
用数学归纳法证明.

当 \(n=1\) 时，等式左端为 \(1^3=1\)，右端为
\[
\left[\frac{1\times\left(1+1\right)}{2}\right]^2=1
\]
所以等式对 \(n=1\) 成立.

假设当 \(n=k\) 时等式成立，即
\[
1^3+2^3+3^3+\cdots+k^3=\left[\frac{k\left(k+1\right)}{2}\right]^2
\]
则
\[
\begin{gathered}
1^3+2^3+3^3+\cdots+k^3+\left(k+1\right)^3\\
=\left[\frac{k\left(k+1\right)}{2}\right]^2+\left(k+1\right)^3=\frac{k^2\left(k+1\right)^2+4\left(k+1\right)^3}{4}\\
=\frac{\left(k+1\right)^2\left(k^2+4k+4\right)}{4}=\frac{\left(k+1\right)^2\left(k+2\right)^2}{4}=\left[\frac{\left(k+1\right)\left(k+2\right)}{2}\right]^2
\end{gathered}
\]
这正是把原等式中的 \(n\) 换成 \(k+1\) 后的右端，因此等式对 \(n=k+1\) 也成立. 由数学归纳法，对任意自然数 \(n\)，都有
\[
1^3+2^3+3^3+\cdots+n^3=\left[\frac{n\left(n+1\right)}{2}\right]^2
\]
故原等式成立.
\end{solution}

\begin{problem}
计算定积分的值：\(\displaystyle\int_0^{\frac{\pi}{4}}\cos2x\,\mathrm{d}x\).
\end{problem}

\begin{solution}
令 \(u=2x\)，则 \(\mathrm{d}u=2\,\mathrm{d}x\)，即 \(\mathrm{d}x=\frac{1}{2}\mathrm{d}u\). 当 \(x=0\) 时 \(u=0\)，当 \(x=\frac{\pi}{4}\) 时 \(u=\frac{\pi}{2}\). 因而
\[
\int_0^{\frac{\pi}{4}}\cos2x\,\mathrm{d}x=\frac{1}{2}\int_0^{\frac{\pi}{2}}\cos u\,\mathrm{d}u=\frac{1}{2}\left[\sin u\right]_0^{\frac{\pi}{2}}=\frac{1}{2}\left(\sin\frac{\pi}{2}-\sin0\right)=\frac{1}{2}
\]
故定积分的值为 \(\frac{1}{2}\).
\end{solution}
